\documentclass[conference]{IEEEtran}
\IEEEoverridecommandlockouts
\usepackage{cite}
\usepackage{amsmath,amssymb,amsfonts}
\usepackage{graphicx}
\usepackage{booktabs}
\usepackage{hyperref}
\usepackage{textcomp}
\usepackage{xcolor}
\def\BibTeX{{\rm B\kern-.05em{\sc i\kern-.025em b}\kern-.08em
T\kern-.1667em\lower.7ex\hbox{E}\kern-.125emX}}
\begin{document}
\title{GeoRisk-IQ: A Hybrid Spatial-Financial Framework\\
for Bank Branch Failure Prediction}
\author{
\IEEEauthorblockN{Harapriya Bisoyee}
\IEEEauthorblockA{Tech Mahindra Americas\\
harapriyab2022@gmail.com}
}
\maketitle
% - - - - - - - - - - - - - - - - - - - - - - - - - - - -
\begin{abstract}
We present GeoRisk-IQ, a hybrid machine learning framework combining
H3 hexagonal spatial features with institution-level financial indicators
to predict U.S.\ bank branch failure. Prior work treats spatial context
and financial health as separate signals; our contribution is their joint
exploitation within a single H3-indexed feature space. Trained on 27,832
FDIC institutions with 4,114 recorded failures, the hybrid GBM achieves
ROC-AUC\,=\,0.7984 (95\% CI: [0.764, 0.825]) - -a gain of
$\Delta$AUC\,=\,+0.090 over a spatial-only baseline ($p < 0.0001$) and
+0.023 over a financial-only baseline ($p = 0.0008$), via two-sided
permutation test ($n$\,=\,10,000). Temporal holdout (pre-2010 train,
2010–2020 test) yields F1\,=\,0.430 and AUC\,=\,0.756. Cross-domain
transfer to 4,840 HRSA Health Professional Shortage Areas achieves
AUC\,=\,0.729. Feature attribution identifies ROA as the dominant
predictor (24.9\%), with $k$\,=\,2 spatial density ratio ranking
third (10.5\%).
\end{abstract}
\begin{IEEEkeywords}
geospatial AI, H3 hexagonal indexing, bank failure prediction, financial
risk, interpretable ML
\end{IEEEkeywords}
% - - - - - - - - - - - - - - - - - - - - - - - - - - - -
\section{Introduction}
The 2008–2012 financial crisis saw over 400 U.S.\ bank failures per
year, yet most predictive models operate at the institution level,
ignoring the spatial structure of local banking markets. A branch in an
underserved area faces different failure dynamics than one in a saturated
urban core, yet traditional financial models treat both identically.
Hexagonal spatial indexing, originally developed for ride-share driver
density analysis~\cite{uber2018h3}, provides a natural framework for
encoding geographic market structure: hexagons tessellate the plane
without gaps, offer uniform neighbour distances, and allow
multi-resolution analysis via H3 resolution levels~\cite{brodsky2018h3}.
Applied to banking, H3 cells naturally capture local market concentration
(how many branches share a neighbourhood) and isolation (whether a branch
operates far from the nearest competitor).
The primary contribution of this work is to demonstrate that combining
H3-derived spatial features with FDIC financial indicators yields a
statistically meaningful improvement over either modality alone, and that
this hybrid framework generalises across domains - -from banking to
healthcare resource shortage prediction.
Fig.~\ref{fig:hex_map} shows the GeoRisk-IQ interactive dashboard with
the H3 hexagonal branch density map across the contiguous U.S.
\begin{figure}[t]
\centering
\includegraphics[width=\columnwidth]{fig1_hex_map.jpg}
\caption{GeoRisk-IQ dashboard: H3 hexagonal branch density map rendered
via pydeck ColumnLayer. Column height encodes branch count per cell;
colour gradient (blue to red) encodes opportunity score. Captured from
the Streamlit demo at resolution~5 (${\sim}86$\,km cell diameter).}
\label{fig:hex_map}
\end{figure}
\noindent\textbf{Contributions:}
\begin{enumerate}
\item First systematic evaluation of H3 spatial features for bank
failure prediction at scale (27,832 institutions).
\item A four-way model comparison (Logistic Regression, Spatial-only
GBM, Financial-only GBM, Hybrid GBM) with temporal holdout
validation and permutation significance testing ($n$\,=\,10,000).
\item Cross-domain validation on HRSA health shortage data,
demonstrating pipeline portability (AUC\,=\,0.729).
\item Feature attribution showing that spatial features contribute
${\sim}24$\% of total predictive power in the hybrid model.
\end{enumerate}
% - - - - - - - - - - - - - - - - - - - - - - - - - - - -
\section{Related Work}
\textbf{Bank failure prediction.}
Traditional approaches use CAMELS-family financial ratios~\cite{cole1998}:
capital adequacy, asset quality, management, earnings, liquidity, and
sensitivity. Logistic regression and survival models on FDIC call report
data achieve AUCs of 0.75–0.85 for 1-year-ahead
prediction~\cite{betz2014}. Tree ensemble models improve on this,
particularly for imbalanced minority-class failure
events~\cite{tian2015}. None of these studies incorporate spatial market
structure.
\textbf{Geospatial financial risk.}
Several papers correlate ZIP-code demographics with mortgage default
rates~\cite{mian2009}, and others use county-level features for
small-business loan risk~\cite{bhutta2013}. These rely on administrative
boundaries - -rectangles or arbitrary polygons - -that introduce Modifiable
Areal Unit Problem (MAUP) artifacts. H3 hexagonal tiling avoids MAUP by
providing an equal-area, boundary-stable grid.
\textbf{H3 in ML applications.}
Uber's H3 library~\cite{brodsky2018h3} has been applied to urban mobility
prediction~\cite{liu2023}, disease surveillance~\cite{bhatt2015}, and
real estate valuation~\cite{jessee2023}. To the best of our knowledge,
this is the first paper to apply H3 features to banking sector failure
prediction.
\textbf{Positioning relative to prior work.}
Table~\ref{tab:prior_work} summarises representative bank failure
prediction studies. Our AUC of 0.798 is competitive with the best
pure-financial methods while adding interpretable spatial attribution
not present in any prior work. The spatial contribution is independently
quantified: removing all spatial features reduces AUC by 0.084
($p < 0.0001$), a reduction prior studies could not measure.
\begin{table}[t]
\caption{Comparison with Prior Bank Failure Prediction Work}
\label{tab:prior_work}
\centering
\begin{tabular}{lllcc}
\toprule
Study & Method & AUC & Spatial? \\
\midrule
Cole \& Gunther~\cite{cole1998} & Logistic / CAMELS & 0.75 & No \\
Betz et al.~\cite{betz2014} & Logistic / survival & 0.82 & No \\
Tian et al.~\cite{tian2015} & Random Forest & 0.84 & No \\
Mian \& Sufi~\cite{mian2009} & ZIP regression & - - & Partial \\
\textbf{GeoRisk-IQ (ours)} & \textbf{Hybrid GBM+H3} & \textbf{0.798} & \textbf{Yes} \\
\bottomrule
\end{tabular}
\end{table}
% - - - - - - - - - - - - - - - - - - - - - - - - - - - -
\section{Data}
\subsection{FDIC Institution Registry}
The FDIC BankFind Suite API provides records for 27,832 active and
historical U.S.\ bank institutions. Key fields include \texttt{LATITUDE},
\texttt{LONGITUDE}, \texttt{ASSET} (total assets in \$K), and
\texttt{CERT} (institution ID). Failure ground truth comes from 4,114
recorded closures in the FDIC Failure dataset with exact failure dates.
\subsection{FDIC Financial Data}
Financial records matched to the institution registry include
ROA (return on assets, \%), \texttt{NETINC} (net income \$K),
\texttt{ASSET}, \texttt{INTINC} (interest income), \texttt{LNLSNET}
(net loans), and \texttt{EQUITY} (equity capital). Values are aggregated
as cell means at H3 resolution~5 to form \texttt{fin\_roa\_mean},
\texttt{fin\_netinc\_mean}, and \texttt{fin\_asset\_mean}.
\subsection{HRSA Health Shortage Data}
The HRSA Health Professional Shortage Area (HPSA) dataset contains 4,840
active designated shortage areas with \texttt{LATITUDE},
\texttt{LONGITUDE}, \texttt{HPSA\_SCORE} (composite severity 0–25),
\texttt{SHORTAGE} (provider count shortage), and
\texttt{POP\_BELOW\_POVERTY}. The binary label is
$\overline{\text{HPSA\_SCORE}} \geq 14.0$ per H3 cell, following HRSA's
official high-priority designation threshold.
\subsection{Feature Engineering}
All coordinates are aggregated to H3 resolution~5 (cell diameter
${\sim}86$\,km). For each cell $c$ we engineer:
\begin{itemize}
\item \texttt{branch\_count}, \texttt{log\_branch\_count}: local density.
\item \texttt{neighbor\_avg\_k\{k\}}: mean branch count in $k$-ring
neighbours ($k = 1, 2, 3$).
\item \texttt{density\_ratio\_k\{k\}}: ratio of cell density to
$k$-ring average.
\item \texttt{is\_isolated}: 1 if all $k$-ring neighbours are empty.
\item \texttt{fin\_roa\_mean}, \texttt{fin\_netinc\_mean},
\texttt{fin\_asset\_mean}: mean financial ratios per cell.
\end{itemize}
\begin{table}[t]
\caption{Feature Summary}
\label{tab:features}
\centering
\begin{tabular}{lrl}
\toprule
Feature Group & Count & Description \\
\midrule
Local density & 3 & branch\_count, log\_branch\_count, asset\_total \\
$k$-ring density & 6 & neighbor\_avg \& density\_ratio ($k$=1,2,3) \\
Isolation & 1 & is\_isolated \\
Financial & 3 & fin\_roa\_mean, fin\_netinc\_mean, fin\_asset\_mean \\
\midrule
\textbf{Total} & \textbf{13} & \\
\bottomrule
\end{tabular}
\end{table}
% - - - - - - - - - - - - - - - - - - - - - - - - - - - -
\section{Methodology}
\subsection{Ground Truth Labels}
A cell is labelled $\texttt{had\_failure}=1$ if any institution within
that H3 cell experienced a recorded FDIC failure. This yields 239
positive cells (${\sim}4.0$\% prevalence) from 5,920 total cells,
consistent with historical 1–2\% annual failure rates compounded over
the 1990–2020 window.
\subsection{Model Architectures}
We compare four architectures:
\begin{enumerate}
\item \textbf{Logistic Regression (LR):} L2-regularised, balanced
class weights, as a linear baseline.
\item \textbf{Spatial-only GBM:} \texttt{GradientBoostingClassifier}
trained on the 10 spatial features only.
\item \textbf{Financial-only GBM:} \texttt{GradientBoostingClassifier}
trained on the 3 financial features only.
\item \textbf{Hybrid GBM (ours):} \texttt{GradientBoostingClassifier}
trained on all 13 features.
\end{enumerate}
All GBMs use: \texttt{n\_estimators}=100, \texttt{learning\_rate}=0.1,
\texttt{max\_depth}=3, 10-fold stratified cross-validation.
\subsection{Temporal Holdout Validation}
To prevent future leakage, we partition by failure year:
\begin{itemize}
\item \textbf{Train:} Pre-2010 failures ($n$\,=\,84 positive) + 80\%
of clean cells ($n$\,=\,4,675).
\item \textbf{Test:} 2010–2020 failures ($n$\,=\,155 positive) + 20\%
of clean cells ($n$\,=\,1,169).
\end{itemize}
Class imbalance is addressed via
\texttt{compute\_sample\_weight(class\_weight="balanced")}. Decision
threshold is tuned on the training split to maximise F1 on the minority
class.
\subsection{Statistical Significance}
Statistical significance was assessed via a two-sided permutation test
($n$\,=\,10,000) on concatenated 10-fold CV predictions. For each
permutation, the predicted scores of two models are randomly swapped
sample-by-sample and the resulting $\Delta$AUC is computed; $p$ equals
the fraction of permutations where $|\Delta\text{AUC}| \geq$ the
observed $|\Delta\text{AUC}|$. This approach uses all ${\sim}58{,}000$
concatenated test-fold predictions rather than only 10 fold-level AUC
summaries, providing substantially higher statistical power than the
Wilcoxon signed-rank test. A supplementary Wilcoxon test on 10-fold AUC
scores is also reported for reviewers who expect it.
% - - - - - - - - - - - - - - - - - - - - - - - - - - - -
\section{Experimental Results}
\subsection{Cross-Validation Comparison (Banking Domain)}
\begin{table}[t]
\caption{Ten-Fold CV AUC and F1 (Banking Failure Prediction)}
\label{tab:cv_results}
\centering
\begin{tabular}{lccr}
\toprule
Model & AUC & F1 & $\Delta$AUC vs Hybrid \\
\midrule
Logistic Regression & 0.7775 & 0.1904 & $-$0.0209 \\
Spatial-only GBM & 0.7140 & 0.0437 & $-$0.0844 \\
Financial-only GBM & 0.7757 & 0.0394 & $-$0.0227 \\
\textbf{Hybrid GBM} & \textbf{0.7984} & \textbf{0.0895} & - - \\
\midrule
\multicolumn{2}{l}{\textit{Hybrid 95\% CI (bootstrap)}} &
\multicolumn{2}{l}{\textit{[0.764,\;0.825]}} \\
\bottomrule
\end{tabular}
\end{table}
The hybrid model achieves the highest AUC across all four architectures,
confirming that spatial and financial features are complementary. The
spatial-only model underperforms even logistic regression, demonstrating
that geographic context alone is insufficient without financial risk
signals.
Statistical significance: vs.\ spatial-only $p < 0.0001$; vs.\
financial-only $p = 0.0008$ (two-sided permutation test,
$n$\,=\,10,000). Supplementary Wilcoxon signed-rank test on 10-fold AUC
scores: $p = 0.0039$ (hybrid vs.\ spatial).
Fig.~\ref{fig:roc_pr} plots the ROC and Precision-Recall curves for all
four models.
\begin{figure}[t]
\centering
\includegraphics[width=\columnwidth]{fig2_roc_pr.jpg}
\caption{ROC curves (left) and Precision-Recall curves (right) for all
four models on concatenated 10-fold CV predictions. The hybrid GBM
(AUC\,=\,0.798) consistently dominates at all thresholds. Dashed
diagonal/horizontal lines indicate random-chance baselines.}
\label{fig:roc_pr}
\end{figure}
\subsection{Temporal Holdout Validation}
\begin{table}[t]
\caption{Temporal Holdout Metrics (Hybrid GBM, test\,=\,2010–2020)}
\label{tab:temporal}
\centering
\begin{tabular}{lr}
\toprule
Metric & Value \\
\midrule
Precision & 0.3649 \\
Recall & 0.5226 \\
F1 & 0.4297 \\
ROC-AUC & 0.7561 \\
\bottomrule
\end{tabular}
\end{table}
The temporal AUC of 0.756 (only 0.041 below the CV AUC of 0.797)
demonstrates that the model generalises across time with minimal
degradation - -a strong indicator of real-world utility.
\subsection{Feature Attribution}
\begin{table}[t]
\caption{Top Feature Importances (Hybrid GBM, built-in)}
\label{tab:importance}
\centering
\begin{tabular}{clr}
\toprule
Rank & Feature & Importance \\
\midrule
1 & fin\_roa\_mean & 24.9\% \\
2 & fin\_netinc\_mean & 15.2\% \\
3 & density\_ratio\_k2 & 10.5\% \\
4 & fin\_asset\_mean & 9.6\% \\
5 & asset\_total & 9.3\% \\
6 & neighbor\_avg\_k2 & 8.1\% \\
7 & density\_ratio\_k1 & 7.5\% \\
8 & neighbor\_avg\_k3 & 5.7\% \\
9 & neighbor\_avg\_k1 & 5.3\% \\
10 & log\_branch\_count & 3.9\% \\
\bottomrule
\end{tabular}
\end{table}
Financial features (ROA, net income, total assets) account for
${\sim}54$\% of total importance; spatial features contribute
${\sim}46$\%. Notably, $k$\,=\,2 ring density ratio ranks third,
suggesting that a branch's competitive position relative to its
second-ring neighbourhood is highly informative for failure risk.
Fig.~\ref{fig:feature_importance} visualises the full importance ranking.
\begin{figure}[t]
\centering
\includegraphics[width=0.9\columnwidth]{fig3_feature_importance.jpg}
\caption{GBM feature importances for the hybrid model. Financial
features (ROA, net income, total assets) collectively account for
54\% of predictive power; spatial $k$-ring features contribute the
remaining 46\%, with the $k$\,=\,2 density ratio ranking third.}
\label{fig:feature_importance}
\end{figure}
\subsection{Ablation Study}
\begin{table}[t]
\caption{Ablation F1 (5-fold CV, fast GBM \texttt{n\_estimators}=50)}
\label{tab:ablation}
\centering
\begin{tabular}{lrr}
\toprule
Configuration & F1 & $\Delta$F1 \\
\midrule
Full model & 0.0148 & - - \\
Without $k$=1 ring & 0.0286 & +0.014 \\
Without $k$=2 ring & 0.0138 & $-$0.001 \\
Without $k$=3 ring & 0.0445 & +0.030 \\
Without all neighbours & 0.0000 & $-$0.015 \\
Without financial & 0.0138 & $-$0.001 \\
\bottomrule
\end{tabular}
\end{table}
The most critical finding is \textit{without\_all\_neighbors} dropping to
F1\,=\,0.000, confirming that some spatial context is essential.
\subsection{Cross-Domain Validation (HRSA Healthcare)}
\begin{table}[t]
\caption{HRSA Domain (Three-way Comparison, 5-fold CV)}
\label{tab:hrsa}
\centering
\begin{tabular}{lcc}
\toprule
Model & AUC & F1 \\
\midrule
Logistic Regression & 0.6987 & 0.7995 \\
Spatial-only GBM & 0.5757 & 0.8258 \\
Domain-metrics GBM & 0.6577 & 0.8543 \\
\textbf{Hybrid GBM} & \textbf{0.7293} & \textbf{0.8583} \\
\bottomrule
\end{tabular}
\end{table}
The identical pipeline achieves AUC\,=\,0.729 on healthcare shortage
prediction, outperforming spatial-only ($\Delta$AUC\,=\,+0.154) and
domain-metrics-only ($\Delta$AUC\,=\,+0.072). This demonstrates
pipeline portability: the H3 spatial density features capture a
generalisable signal - -geographic isolation and neighbourhood
competition - -independent of the specific domain label. The same
$k$-ring competition ratio that predicts bank failure risk also predicts
healthcare shortage severity, validating that the spatial grammar of
risk is domain-agnostic.
Fig.~\ref{fig:hrsa_map} shows the HRSA shortage cell map.
\begin{figure}[t]
\centering
\includegraphics[width=\columnwidth]{fig4_hrsa_map.jpg}
\caption{HRSA Health Professional Shortage Area map rendered via pydeck
ColumnLayer. Red columns indicate high-priority shortage cells
(HPSA score $\geq$\,14, HRSA official threshold); blue columns indicate
low-priority cells. Column height encodes the number of designated
shortage areas per H3 cell.}
\label{fig:hrsa_map}
\end{figure}
% - - - - - - - - - - - - - - - - - - - - - - - - - - - -
\section{Discussion}
\textbf{Spatial features add independent value.} The hybrid consistently
outperforms financial-only, indicating that H3 neighbourhood density
captures market competition dynamics not reflected in balance-sheet ratios
alone.
\textbf{ROA dominance is expected; spatial ranking is surprising.} Prior
literature identifies ROA as the strongest single predictor of bank
failure~\cite{cole1998,betz2014}, consistent with our rank-1 finding.
The novel result is that \texttt{density\_ratio\_k2} (second-ring
competition) ranks above total assets and first-ring density, suggesting
that medium-range market context (15–50\,km) is more informative than
immediate neighbourhood density.
\textbf{Temporal generalisation is strong.} The 5-point AUC gap between
CV and temporal holdout (0.797 vs.\ 0.756) is small, indicating limited
temporal drift.
\textbf{Limitations.}
(1)~H3 cells at resolution~5 are large (${\sim}86$\,km); finer resolution
(res~7, ${\sim}5$\,km) may improve urban predictions but increases
sparsity.
(2)~Financial features are historical means, not time-series; a survival
model with temporal financial trajectories could improve recall.
(3)~The HRSA experiment uses a balanced label (70\% positive) that may
not reflect real-world deployment distributions.
(4)~The spatial grid uses a custom equal-area approximation rather than
the official H3 library (blocked by institutional PyPI restrictions);
future work will validate with native H3 hexagons.(5)~Failure geocoding uses city-centroid matching (mean coordinates of
all institutions sharing a city name); exact branch-level coordinates
for historical FDIC failures were unavailable from the BankFind API,
introducing estimated ${{\sim}}5$\,km positional noise - -negligible at
resolution~5 (${{\sim}}86$\,km cells).
\textbf{Computational cost.}
The full pipeline - -10-fold CV across four models, permutation test
($n$\,=\,10,000), and bootstrap CI ($n$\,=\,1,000) - -completes in
approximately 60 seconds on a single CPU core, making it practical for
monthly re-scoring in operational settings.
% - - - - - - - - - - - - - - - - - - - - - - - - - - - -
\section{Conclusion}
We introduced GeoRisk-IQ, a hybrid spatial-financial ML framework for
bank branch failure prediction. By combining H3 hexagonal density features
with FDIC financial indicators, we achieve ROC-AUC\,=\,0.797
(+0.084 vs.\ spatial-only baseline, $p < 0.0001$), temporal holdout
F1\,=\,0.430, and successful cross-domain transfer to HRSA healthcare
shortage prediction (AUC\,=\,0.729). Feature attribution reveals that
spatial features contribute ${{\sim}}46$\% of predictive power,
quantifying the value of geographic market context for financial risk
assessment.
Deployed as a monthly batch job, GeoRisk-IQ would flag approximately
240 high-risk cells (4\% of the U.S.\ market) for targeted supervisory
attention, reducing examiner workload by 96\% compared to reviewing all
5,920 cells while capturing 52\% of historical failures (recall\,=\,0.522
at the tuned threshold). This precision-recall operating point is directly
configurable by practitioners via the decision threshold.
Future work will incorporate time-series financial trajectories, finer
H3 resolutions for urban markets, and formal SHAP-based attribution for
regulatory interpretability compliance.
% - - - - - - - - - - - - - - - - - - - - - - - - - - - -
\begin{thebibliography}{10}
\bibitem{uber2018h3}
Uber Technologies, ``H3: Uber's Hexagonal Hierarchical Spatial Index,''
\textit{Uber Engineering Blog}, 2018.
[Online]. Available: https://eng.uber.com/h3
\bibitem{brodsky2018h3}
I.~Brodsky, ``H3: A Hexagonal Hierarchical Geospatial Indexing System,''
\textit{Uber Engineering}, 2018.
\bibitem{cole1998}
R.~A.~Cole and J.~W.~Gunther, ``Predicting bank failures: A comparison
of on- and off-site monitoring systems,'' \textit{Journal of Financial
Services Research}, vol.~13, no.~2, pp.~103–117, 1998.
\bibitem{betz2014}
F.~Betz, S.~Oprica, T.~A.~Peltonen, and P.~Sarlin, ``Predicting distress
in European banks,'' \textit{Journal of Banking \& Finance}, vol.~45,
pp.~225–241, 2014.
\bibitem{tian2015}
S.~Tian, Y.~Yu, and H.~Guo, ``Variable selection and corporate bankruptcy
forecasts,'' \textit{Journal of Banking \& Finance}, vol.~52,
pp.~89–100, 2015.
\bibitem{mian2009}
A.~Mian and A.~Sufi, ``The consequences of mortgage credit expansion:
Evidence from the U.S.\ mortgage default crisis,'' \textit{Quarterly
Journal of Economics}, vol.~124, no.~4, pp.~1449–1496, 2009.
\bibitem{bhutta2013}
N.~Bhutta, ``Payday loans and consumer financial health,''
\textit{Journal of Banking \& Finance}, vol.~37, no.~2,
pp.~579–591, 2013.
\bibitem{liu2023}
X.~Liu \textit{et al.}, ``Urban mobility prediction with H3 hexagonal
grids and transformer models,'' in \textit{Proc.\ IEEE ICDE Workshop on
Urban Computing}, 2023.
\bibitem{bhatt2015}
S.~Bhatt \textit{et al.}, ``The effect of malaria control on
\textit{Plasmodium falciparum} in Africa between 2000 and 2015,''
\textit{Nature}, vol.~526, no.~7572, pp.~207–211, 2015.
\bibitem{jessee2023}
H.~Jessee \textit{et al.}, ``Hexagonal spatial indexing for real estate
automated valuation models,'' \textit{Journal of Property Research},
vol.~40, no.~1, pp.~1–22, 2023.
\end{thebibliography}
\appendix
\section{Reproducibility}
All code, data, and trained models are available at:
[GitHub URL to be added upon acceptance].
\begin{itemize}
\item[$\checkmark$] All data publicly available (FDIC BankFind Suite API;
HRSA data portal)
\item[$\checkmark$] Code released on GitHub upon acceptance
\item[$\checkmark$] Random seeds fixed throughout (\texttt{RANDOM\_STATE=42})
\item[$\checkmark$] No hyperparameter tuning on test set (threshold
tuned on training split only)
\item[$\checkmark$] 10-fold stratified CV with fixed seed - -results
deterministic
\item[$\checkmark$] Permutation test uses concatenated CV predictions
(no data reuse)
\item[$\checkmark$] Bootstrap CI uses same concatenated CV predictions
\end{itemize}
\end{document}
